\documentclass[oneside]{book}
\usepackage{amsmath, amssymb, amsthm}
\usepackage{environ}


% 布局
\usepackage{geometry}
\geometry{
  a4paper,           % 纸张：A4
  left=25mm,         % 左页边距 20mm
  right=25mm,        % 右页边距 20mm
  top=30mm,          % 上页边距 25mm
  bottom=30mm,       % 下页边距 25mm
  includefoot,       % 关键：将页脚区域计入总版心高度
%   includehead        % 关键：将页眉区域计入总版心高度
}

% 颜色支持（可选，用于答案高亮）
\usepackage{xcolor}

% 超链接（用于PDF输出）
\usepackage[hidelinks]{hyperref}
\usepackage{bookmark}

\usepackage{ctex}

\usepackage{tikz}

\usepackage{tkz-euclide} %为几何画图设计的宏包
\usetikzlibrary{calc} % 坐标计算功能
\usetikzlibrary{intersections} % 交点计算功能
\usetikzlibrary{arrows.meta} %箭头
\tikzset{>={Stealth[scale=1.4]}} %燕尾箭头

\usepackage[most]{tcolorbox}
\tcbuselibrary{breakable}

\usepackage{titlesec}

% 修改chapter格式
\titleformat{\chapter}
  {\centering\LARGE\bfseries}  % 格式
  {第\,\thechapter\,章}        % 标签
  {1em}                       % 间距
  {}                          % 标题前的内容


% 控制答案显隐的开关
\newif\ifshowanswers
\showanswerstrue % 显示答案（编译时注释掉这行即可隐藏）


% 答案环境定义 - 直接集成彩色框
\NewEnviron{answer}{
    \ifshowanswers
        \begin{tcolorbox}[
            colback=red!5!white,
            colframe=red!75!black,
            fonttitle=\bfseries,
            title=答案,
            breakable,
        ]
        \BODY
        \end{tcolorbox}
    \fi
}

% 一些颜色 violet, olive
\newtcolorbox[auto counter,number within=chapter]{definition}[1][]{%
colback=violet!5!white,
colframe=violet!75!black,
coltitle=white,
colbacktitle=violet!75!black,
fonttitle=\bfseries,
title=定义~\thetcbcounter~#1,
breakable,}

\newtcolorbox[auto counter,number within=chapter]{example}[1][]{%
colback=green!5!white,
colframe=green!75!black,
% sharp corners,
% rounded corners=uphill,
fonttitle=\bfseries,
title=例题~\thetcbcounter~#1,
breakable,}

\newtcolorbox[auto counter,number within=chapter]{sectionexercise}[1][]{%
colback=cyan!5!white,
colframe=cyan!75!black,
fonttitle=\bfseries,
title=章节练习~\thetcbcounter~#1,
breakable,}






\begin{document}

\title{数学讲义}
\author{QQ群 827687129}
\date{\today}

\maketitle

\tableofcontents


\chapter{微积分基础}

本章将介绍微积分的基本概念，包括极限、导数和积分，并通过例题和练习加深理解。

\newpage

\section{极限与连续性}

\subsection{极限的概念}

\begin{definition}[极限的$\varepsilon$-$\delta$定义]
设函数$f(x)$在点$x_0$的某个去心邻域内有定义，如果存在常数$A$，对于任意给定的正数$\varepsilon$（无论它多么小），总存在正数$\delta$，使得当$0<|x-x_0|<\delta$时，对应的函数值$f(x)$都满足不等式$|f(x)-A|<\varepsilon$，那么常数$A$就叫做函数$f(x)$当$x \to x_0$时的极限，记作：
\[
\lim_{x \to x_0} f(x) = A
\]
\end{definition}

\begin{example}[基本极限计算]
计算极限：$\displaystyle \lim_{x \to 0} \frac{\sin x}{x}$。
\end{example}

\begin{answer}
使用重要极限公式：
\[
\lim_{x \to 0} \frac{\sin x}{x} = 1
\]
可以通过洛必达法则证明：
\[
\lim_{x \to 0} \frac{\sin x}{x} = \lim_{x \to 0} \frac{\cos x}{1} = \cos 0 = 1
\]
或者使用夹逼定理证明。
\end{answer}

\subsection{连续性}

\begin{definition}[函数连续性]
设函数$f(x)$在点$x_0$的某个邻域内有定义，如果
\[
\lim_{x \to x_0} f(x) = f(x_0)
\]
则称函数$f(x)$在点$x_0$处连续。
\end{definition}

\begin{example}[判断函数连续性]
判断函数$f(x) = \begin{cases} 
x^2 + 1, & x \geq 0 \\
x - 1, & x < 0 
\end{cases}$在$x=0$处的连续性。
\end{example}

\begin{answer}
计算左右极限：
\[
\lim_{x \to 0^-} f(x) = \lim_{x \to 0^-} (x-1) = -1
\]
\[
\lim_{x \to 0^+} f(x) = \lim_{x \to 0^+} (x^2+1) = 1
\]
由于左极限不等于右极限，所以函数在$x=0$处不连续。
\end{answer}

\begin{sectionexercise}[极限与连续性]
\begin{enumerate}

    \item 用 $\varepsilon$-$\delta$ 语言证明：$\displaystyle \lim_{x \to 2} (3x-1) = 5$。
    
    \item 计算下列极限：
    \begin{enumerate}
        \item $\displaystyle \lim_{x \to 0} \frac{\tan 3x}{x}$
        \item $\displaystyle \lim_{x \to 4} \frac{\sqrt{x} - 2}{x - 4}$
        \item $\displaystyle \lim_{x \to \infty} \frac{3x^2 + 2x + 1}{2x^2 - x + 3}$
    \end{enumerate}

    \item 讨论函数 $f(x) = \begin{cases}
        \dfrac{\sin 2x}{x}, & x \neq 0 \\
        2, & x = 0
    \end{cases}$ 在 $x=0$ 处的连续性。
    
    

    \item 设 $f(x) = \begin{cases}
        x^2 + ax + b, & x \leq 1 \\
        2x + 3, & x > 1
    \end{cases}$，问常数 $a, b$ 取何值时，函数 $f(x)$ 在 $x=1$ 处连续？
    
    \item 判断下列命题是否正确，并说明理由：
    \begin{enumerate}
        \item 若 $\displaystyle \lim_{x \to x_0} f(x)$ 存在，则 $f(x)$ 在 $x_0$ 处连续。
        \item 若 $f(x)$ 在 $x_0$ 处连续，则 $\displaystyle \lim_{x \to x_0} f(x)$ 存在。
        \item 若 $f(x)$ 在 $x_0$ 处不连续，则 $\displaystyle \lim_{x \to x_0} f(x)$ 不存在。
    \end{enumerate}

\end{enumerate}
\end{sectionexercise}

\begin{answer}
\begin{enumerate}
    \item 证明：对于任意给定的 $\varepsilon > 0$，要找到 $\delta > 0$，使得当 $0 < |x-2| < \delta$ 时，有 $|(3x-1)-5| < \varepsilon$。
    
    因为 $|(3x-1)-5| = |3x-6| = 3|x-2|$，所以要使 $3|x-2| < \varepsilon$，只需取 $\delta = \varepsilon/3$。
    
    于是，对于任意 $\varepsilon > 0$，取 $\delta = \varepsilon/3$，当 $0 < |x-2| < \delta$ 时，有 $|(3x-1)-5| = 3|x-2| < 3 \cdot (\varepsilon/3) = \varepsilon$。
    
    因此，$\displaystyle \lim_{x \to 2} (3x-1) = 5$。
    
    \item 计算下列极限：
    \begin{enumerate}
        \item $\displaystyle \lim_{x \to 0} \frac{\tan 3x}{x} = \lim_{x \to 0} \frac{\sin 3x}{x \cos 3x} = \lim_{x \to 0} \frac{3\sin 3x}{3x} \cdot \frac{1}{\cos 3x} = 3 \times 1 \times 1 = 3$
        
        \item $\displaystyle \lim_{x \to 4} \frac{\sqrt{x} - 2}{x - 4} = \lim_{x \to 4} \frac{(\sqrt{x} - 2)(\sqrt{x} + 2)}{(x-4)(\sqrt{x} + 2)} = \lim_{x \to 4} \frac{x-4}{(x-4)(\sqrt{x} + 2)} = \lim_{x \to 4} \frac{1}{\sqrt{x} + 2} = \frac{1}{2+2} = \frac{1}{4}$
        
        \item $\displaystyle \lim_{x \to \infty} \frac{3x^2 + 2x + 1}{2x^2 - x + 3} = \lim_{x \to \infty} \frac{3 + \frac{2}{x} + \frac{1}{x^2}}{2 - \frac{1}{x} + \frac{3}{x^2}} = \frac{3 + 0 + 0}{2 - 0 + 0} = \frac{3}{2}$
    \end{enumerate}
    
    \item 讨论函数在 $x=0$ 处的连续性：
    
    计算 $\displaystyle \lim_{x \to 0} f(x) = \lim_{x \to 0} \frac{\sin 2x}{x} = \lim_{x \to 0} \frac{2\sin 2x}{2x} = 2 \times 1 = 2$
    
    而 $f(0) = 2$，所以 $\displaystyle \lim_{x \to 0} f(x) = f(0) = 2$。
    
    因此，函数 $f(x)$ 在 $x=0$ 处连续。
    
    \item 要使函数在 $x=1$ 处连续，需要满足：
    \[
    \lim_{x \to 1^-} f(x) = \lim_{x \to 1^+} f(x) = f(1)
    \]
    
    计算左极限：$\displaystyle \lim_{x \to 1^-} f(x) = \lim_{x \to 1^-} (x^2 + ax + b) = 1^2 + a \cdot 1 + b = 1 + a + b$
    
    计算右极限：$\displaystyle \lim_{x \to 1^+} f(x) = \lim_{x \to 1^+} (2x + 3) = 2 \times 1 + 3 = 5$
    
    函数值：$f(1) = 1^2 + a \cdot 1 + b = 1 + a + b$
    
    由连续性条件得：
    \[
    1 + a + b = 5
    \]
    解得 $a + b = 4$。
    
    因此，当 $a + b = 4$ 时，函数 $f(x)$ 在 $x=1$ 处连续。例如，取 $a = 0, b = 4$ 或 $a = 1, b = 3$ 等。
    
    \item 判断命题：
    \begin{enumerate}
        \item 错误。$\displaystyle \lim_{x \to x_0} f(x)$ 存在是 $f(x)$ 在 $x_0$ 处连续的必要条件，但不是充分条件。还需要满足 $f(x_0)$ 存在且等于极限值。
        
        \item 正确。这是连续性的定义：若 $f(x)$ 在 $x_0$ 处连续，则 $\displaystyle \lim_{x \to x_0} f(x) = f(x_0)$，所以极限存在。
        
        \item 错误。函数在某点不连续，极限仍可能存在。例如：
        \[
        f(x) = \begin{cases}
            x^2, & x \neq 0 \\
            1, & x = 0
        \end{cases}
        \]
        在 $x=0$ 处不连续，但 $\displaystyle \lim_{x \to 0} f(x) = 0$ 存在。
    \end{enumerate}
\end{enumerate}
\end{answer}



\newpage
\section{导数及其应用}

\subsection{导数的定义}

\begin{definition}[导数]
设函数$y=f(x)$在点$x_0$的某个邻域内有定义，当自变量$x$在$x_0$处取得增量$\Delta x$（点$x_0+\Delta x$仍在该邻域内）时，相应地函数取得增量$\Delta y = f(x_0+\Delta x) - f(x_0)$；如果$\Delta y$与$\Delta x$之比当$\Delta x \to 0$时的极限存在，则称函数$y=f(x)$在点$x_0$处可导，并称这个极限为函数$y=f(x)$在点$x_0$处的导数，记为$f'(x_0)$：
\[
f'(x_0) = \lim_{\Delta x \to 0} \frac{\Delta y}{\Delta x} = \lim_{\Delta x \to 0} \frac{f(x_0+\Delta x) - f(x_0)}{\Delta x}
\]
\end{definition}

\begin{example}[使用定义求导]
使用导数定义求函数$f(x)=x^2$在$x=1$处的导数。
\end{example}

\begin{answer}
根据导数定义：
\[
f'(1) = \lim_{\Delta x \to 0} \frac{f(1+\Delta x) - f(1)}{\Delta x}
= \lim_{\Delta x \to 0} \frac{(1+\Delta x)^2 - 1^2}{\Delta x}
= \lim_{\Delta x \to 0} \frac{1 + 2\Delta x + (\Delta x)^2 - 1}{\Delta x}
\]
\[
= \lim_{\Delta x \to 0} \frac{2\Delta x + (\Delta x)^2}{\Delta x}
= \lim_{\Delta x \to 0} (2 + \Delta x) = 2
\]
所以$f'(1)=2$。
\end{answer}

\subsection{导数的基本公式}

\begin{example}[导数计算]
求函数$f(x)=3x^4 - 2x^3 + 5x - 7$的导数。
\end{example}

\begin{answer}
使用幂函数求导公式：
\[
f'(x) = \frac{d}{dx}(3x^4) - \frac{d}{dx}(2x^3) + \frac{d}{dx}(5x) - \frac{d}{dx}(7)
\]
\[
= 3 \cdot 4x^{3} - 2 \cdot 3x^{2} + 5 \cdot 1 - 0
= 12x^3 - 6x^2 + 5
\]
\end{answer}

\newpage
\section{积分学基础}

\subsection{不定积分}

\begin{definition}[不定积分]
设$F(x)$是函数$f(x)$在区间$I$上的一个原函数，则$f(x)$的所有原函数称为$f(x)$在区间$I$上的不定积分，记作：
\[
\int f(x)dx = F(x) + C
\]
其中$C$为任意常数。
\end{definition}

\begin{example}[基本积分计算]
计算不定积分：$\displaystyle \int (3x^2 + 2x + 1)dx$。
\end{example}

\begin{answer}
使用幂函数积分公式：
\[
\int (3x^2 + 2x + 1)dx = \int 3x^2 dx + \int 2x dx + \int 1 dx
\]
\[
= 3 \cdot \frac{x^{3}}{3} + 2 \cdot \frac{x^{2}}{2} + x + C
= x^3 + x^2 + x + C
\]
其中$C$为积分常数。
\end{answer}

\newpage
\subsection{定积分}

\begin{definition}[定积分]
设函数$f(x)$在区间$[a,b]$上有界，在$[a,b]$中任意插入若干个分点：
\[
a = x_0 < x_1 < x_2 < \cdots < x_{n-1} < x_n = b
\]
把区间$[a,b]$分成$n$个小区间，在每个小区间$[x_{i-1}, x_i]$上任取一点$\xi_i$，作和式：
\[
S_n = \sum_{i=1}^{n} f(\xi_i) \Delta x_i
\]
其中$\Delta x_i = x_i - x_{i-1}$。如果当$\lambda = \max\{\Delta x_1, \Delta x_2, \ldots, \Delta x_n\} \to 0$时，和式$S_n$的极限存在，且与区间$[a,b]$的分法及点$\xi_i$的取法无关，则称函数$f(x)$在区间$[a,b]$上可积，并称此极限为$f(x)$在区间$[a,b]$上的定积分，记作：
\[
\int_{a}^{b} f(x)dx = \lim_{\lambda \to 0} \sum_{i=1}^{n} f(\xi_i) \Delta x_i
\]
\end{definition}

\begin{example}[定积分计算]
计算定积分：$\displaystyle \int_{0}^{1} (2x + 1)dx$。
\end{example}

\begin{answer}
首先求原函数：
\[
\int (2x + 1)dx = x^2 + x + C
\]
然后使用牛顿-莱布尼茨公式：
\[
\int_{0}^{1} (2x + 1)dx = \left[ x^2 + x \right]_{0}^{1} = (1^2 + 1) - (0^2 + 0) = 2
\]
\end{answer}




\end{document}